% Generated draft: requires mathematical and source review.
\documentclass[11pt]{article}
\usepackage[margin=25mm]{geometry}
\usepackage{fontspec}
\setmainfont{Latin Modern Roman}
\usepackage{amsmath,amssymb,mathtools,unicode-math}
\setmathfont{Latin Modern Math}
\usepackage{enumitem,needspace,xurl,hyperref}
\setlength{\parindent}{0pt}
\setlength{\parskip}{5pt}
\setlength{\emergencystretch}{3em}
\newcommand{\N}{\mathbb{N}}
\newcommand{\Z}{\mathbb{Z}}
\newcommand{\Q}{\mathbb{Q}}
\newcommand{\R}{\mathbb{R}}
\newcommand{\C}{\mathbb{C}}
\newcommand{\F}{\mathbb{F}}
\newcommand{\A}{\mathbb{A}}
\newcommand{\PP}{\mathbb{P}}
\newcommand{\E}{\mathbb{E}}
\newcommand{\eps}{\varepsilon}
\newcommand{\dd}{\,\mathrm{d}}
\begin{document}
\section*{UnsolvedMath ID: 800979}
\textbf{Source problem code:} catalogue entry
\subsection{Definitions and mathematical statement}
Let \(T^2=\R^2/\Z^2\). Under the usual smooth, globally defined interpretation, let \(g\) be a smooth Riemannian metric on \(T^2\). A polynomial first integral of degree exactly three in the velocities is a smooth function \(F:TT^2\to\R\) of the form
\[
F(x,v)=A_0(x)+A_1(x)(v)+A_2(x)(v,v)+A_3(x)(v,v,v),
\]
where \(A_k\) is a smooth symmetric covariant \(k\)-tensor field and \(A_3\not\equiv0\), such that, for every affinely parametrized \(g\)-geodesic \(\gamma:I\to T^2\),
\[
\frac{\dd}{\dd t}F\bigl(\gamma(t),\dot\gamma(t)\bigr)=0
\qquad(t\in I).
\]
A nonzero Killing vector field is a smooth vector field \(X\not\equiv0\) satisfying \(\mathcal L_Xg=0\); equivalently, with \(\nabla\) the Levi--Civita connection,
\[
g(\nabla_YX,Z)+g(Y,\nabla_ZX)=0
\]
for all smooth vector fields \(Y,Z\).

\textbf{Conjecture 60, under these conventions.} For every such metric \(g\), if its geodesic flow admits a polynomial first integral of degree exactly three in the velocities, then \(g\) admits a nonzero Killing vector field.

\textbf{Unresolved qualifications.} The supplied wording does not specify smoothness, global versus local integrals, or whether the cubic integral must be homogeneous. The formulation above explicitly uses smooth global data and allows lower-degree terms. For the homogeneous reading, require \(A_0=A_1=A_2=0\). These conventions require source review. The word ``nonzero'' is also implicit rather than explicit in the supplied conclusion; without it, the zero vector field would satisfy the conclusion for every metric.
\subsection{Short English statement}
The supplied conjecture says that a Riemannian metric on the two-dimensional torus whose geodesic flow has a cubic polynomial conserved quantity must have a nonzero infinitesimal isometry. The precise regularity, globality, and homogeneity conventions require review.
\subsection{Sources}
\begin{itemize}
\item Local UnsolvedMath registry, supplied complete source record: permanent ID \(800979\), exact source code \texttt{catalogue entry}; upstream revision \texttt{d371ee917ac9647fbb06cd6fc529deec6f3dda82}.
\item Burns and Matveev, \emph{Open Problems and Questions about Geodesics} (2019, as identified by the supplied record), Conjecture 60. Supplied source URL: \url{https://arxiv.org/abs/1308.5417}. The supplied provenance records access on 2026-10-16 and extraction by \texttt{arxiv\_latex}; these are supplied provenance claims, not independently verified access.
\item The source title and background associate the conjecture with the reference label \texttt{[Bo-Ko-Fo1995]}. No corresponding bibliographic entry is supplied, so its authorship and publication details remain unresolved. The registry lists Burns and Matveev as proposers and 2019 as the proposed year; this does not resolve the attribution represented by that reference label.
\item The registry labels the problem ``open''. This is a supplied status label, not an independently verified assessment of the literature.
\end{itemize}
\end{document}
